% Archivo auto-generado por IA y datos del modelo. No editar manualmente.
\newcommand{\ipaUltimoPeriodo}{2026Q1}
\newcommand{\ipaRangoAnos}{2026}
\newcommand{\ipaPeriodoInflexion}{2023T3}
\newcommand{\ipaGlobalScore}{56,4}
\newcommand{\ipaGlobalScoreEsp}{60,4}
\newcommand{\ipaGlobalIncInicio}{+5,3}
\newcommand{\ipaGlobalPctInicio}{+10,5}
\newcommand{\ipaGlobalIncReciente}{+3,6}
\newcommand{\ipaGlobalPctReciente}{+6,8}
\newcommand{\ipaScoreSocInclusivas}{50,6}
\newcommand{\ipaIncSocInclusivasReciente}{-0,4}
\newcommand{\ipaPctIncSocInclusivasReciente}{-0,8}
\newcommand{\ipaPctIncSocInclusivasRecienteEsp}{-28,2}
\newcommand{\ipaDiferenciaEspSocInclusivas}{+8,6}
\newcommand{\ipaScoreEcoAbiertas}{50,5}
\newcommand{\ipaIncEcoAbiertasInicio}{+11,7}
\newcommand{\ipaPctIncEcoAbiertasInicio}{+30,0}
\newcommand{\ipaPctIncEcoAbiertasInicioEsp}{+52,9}
\newcommand{\ipaPctIncEcoAbiertasReciente}{+22,3}
\newcommand{\ipaPctIncEcoAbiertasRecienteEsp}{+42,7}
\newcommand{\ipaMaxEcoAbiertas}{53,0}
\newcommand{\ipaScorePersonas}{70,4}
\newcommand{\ipaIncPersonasInicio}{+9,4}
\newcommand{\ipaPctIncPersonasInicio}{+15,4}
\newcommand{\ipaPctIncPersonasInicioEsp}{+31,8}
\newcommand{\ipaGlobalInterpretation}{\begin{itemize}
    \item Comparativa \textbf{2016Q1 vs 2026Q1}: Incremento absoluto de \textbf{+5,3 puntos} en el nivel de prosperidad.
    \vspace{0.2cm}
    \item Comparativa \textbf{2023T3 vs 2026Q1}: Crecimiento absoluta de \textbf{+3,6 puntos} en la etapa reciente.
    \vspace{0.2cm}
    \item Los motores en el periodo \textbf{2023T3 vs 2026Q1}: \textit{Economías Abiertas} (\textbf{+9,2 puntos}) y \textit{Personas Empoderadas} (\textbf{+0,3 puntos}).
\end{itemize}}
\newcommand{\ipaInclusivasInterpretation}{\begin{itemize}
    \item Comparativa \textbf{2023T3 vs 2026Q1}: Caída absoluta de \textbf{-0,4 puntos} en Andalucía, situando a la comunidad en los \textbf{50,6 puntos}.
    \vspace{0.1cm}
    \item \textbf{Brecha Nacional (2026Q1)}: La puntuación se sitúa a \textbf{+8,6 puntos} del promedio nacional (42,0 puntos).
    \vspace{0.1cm}
    \item Comparativa de variación (\textbf{2023T3 vs 2026Q1}): Andalucía cae \textbf{-0,4 puntos} frente a la variación de \textbf{-16,5 puntos} a nivel nacional.
\end{itemize}}
\newcommand{\ipaEconomiasInterpretation}{\begin{itemize}
    \item Comparativa \textbf{2016Q1 vs 2026Q1}: Crecimiento absoluto de \textbf{+11,7 puntos} en Andalucía frente a \textbf{+22,2 puntos} en España.
    \vspace{0.1cm}
    \item Comparativa reciente (\textbf{2023T3 vs 2026Q1}): El ecosistema andaluz avanza \textbf{+9,2 puntos} frente a \textbf{+19,2 puntos} en España.
    \vspace{0.1cm}
    \item \textbf{Resultado actual (2026Q1)}: Se alcanzan los \textbf{50,5 puntos} (el máximo de la serie histórica se sitúa en \textbf{53,0 puntos}).
\end{itemize}}
\newcommand{\ipaPersonasInterpretation}{\begin{itemize}
    \item Comparativa \textbf{2016Q1 vs 2026Q1}: Incremento absoluto de \textbf{+9,4 puntos} en Andalucía, frente a un aumento de \textbf{+19,7 puntos} a nivel nacional.
    \vspace{0.1cm}
    \item \textbf{Consolidación en el último periodo (2026Q1)}: Puntuación de \textbf{70,4 puntos} en el último trimestre.
\end{itemize}}
